\section{Related Work}


The concept of DTs was introduced in~\cite{grieves2005} and later formalized for CPSs~\cite{glaessgen2012}. DTs are now widely used as evolving software representations of physical assets driven by real-world data streams~\cite{botin2022digital,mihai2022digital,10.1145/3735948.3736156}.


Ontologies and knowledge graphs are widely used in DT architectures for knowledge representation and interoperability~\cite{KARABULUT2024442}. A common approach employs ontologies as schema-level bridges between sensor data and DT models~\cite{10.1016/j.rcim.2023.102649,steinmetz2022methodology}. In contrast, our framework goes beyond structural compatibility and captures behavioral correspondence: semantic alignment requires equivalence of observable behaviors, since twins may conform to the same schema while exhibiting different dynamics. Ontologies are also used as semantic abstractions for cross-domain reasoning, service engineering, and orchestration of DT services~\cite{Boje2024DTSemantics,elhajj2026semantic,10.1145/3652620.3688261,GrangelGonzalez2016,Bader2019SemanticAAS,Arm2021AAS,farias2024dtag}. In these approaches, ontologies support conceptual interoperability between heterogeneous artifacts. In our framework, instead, the ontology acts as a formal semantic ground assigning meaning to PT and DT states and labels, enabling alignment to be defined as behavioral equivalence and guaranteeing preservation of temporal properties across all reachable states.


%Ontologies and knowledge graphs are widely used in DT architectures for knowledge representation, schema integration, and interoperability~\cite{KARABULUT2024442}. A common approach uses ontologies as schema-level bridges between sensor data and DT models~\cite{10.1016/j.rcim.2023.102649,steinmetz2022methodology}. In contrast, our framework goes beyond structural compatibility to capture behavioral correspondence: semantic alignment requires not only compatible data schemas but equivalence of observable behaviors, since twins may conform to the same schema while exhibiting different dynamics.

%Ontologies are also used as higher-level semantic abstractions for cross-domain reasoning, service engineering, and orchestration of DT services~\cite{Boje2024DTSemantics,elhajj2026semantic,10.1145/3652620.3688261}. In these approaches, ontologies primarily support conceptual interoperability between heterogeneous engineering artifacts. Our work differs in that the ontology acts as a formal semantic ground assigning meaning to the states and labels of the PT and DT models themselves. This enables semantic alignment to be defined as a behavioral equivalence relation, yielding preservation guarantees for temporal properties across all reachable states rather than compatibility at the architectural or service level.

PT/DT alignment is often addressed through runtime monitoring and trace analysis. Munoz et al.~\cite{10.1145/3652620.3688267} detect divergence by aligning PT traces with DT behavior and applying anomaly detection, while Moon et al.~\cite{moon2023consistency} verify that observed timed traces conform to a DT model expressed as a Time Colored Petri Net. These approaches rely on syntactic correspondence between traces or executions. In contrast, our framework grounds PT and DT observations in a shared ontology, enabling semantic equivalence even when their observable vocabularies differ. Ontologies are similarly used in~\cite{kamburjan2022twinningbyconstruction,Gil2025} to lift DT runtime states into knowledge graphs for semantic monitoring of PT/DT correspondence.

More generally, runtime monitoring approaches establish alignment operationally and over selected constraints, whereas our framework guarantees semantic consistency entirely at design time over all reachable states and all properties expressible over the aligned fragment. This removes the engineering overhead of implementing and maintaining monitoring infrastructures as systems evolve, where changes would require re-validating and updating all monitors. Our framework further supports ontology evolution without invalidating previously established alignment. Runtime monitoring nevertheless remains complementary for detecting deployment-time deviations~\cite{10.1145/3735948.3736156} or quantifying uncertainty under incomplete environment models~\cite{Woodcock2021}, which lie outside the scope of our work.

Recent work also applies formal methods to DT behavior. Huang et al.~\cite{huang2024formaldt} use the Temporal Logic of Actions to verify synchronization and information-flow properties of DT architectures, while Liu et al.~\cite{liu2023spn} introduce Semantic Petri Nets that combine RDF/OWL knowledge bases with Petri Net models to validate temporal rules over DT executions. These approaches verify properties of the DT itself, whereas our work formalizes semantic alignment between PT and DT as a behavioral equivalence relation under a shared semantic.


Another line of work studies co-evolution of PT and DT models. Model-driven toolchains propagate structural changes from the PT to maintain DT infrastructure~\cite{lehner2021aml4dt,Mertens2024}. These approaches ensure structural correspondence between twins but do not address semantic alignment. Our framework is complementary: it does not prescribe how the DT evolves but provides guarantees that the resulting models remain semantically aligned with respect to domain knowledge.



Taken together, existing work addresses interoperability, runtime monitoring, or verification of DT behavior, but none establishes semantic alignment between PT and DT as a formally verified behavioral equivalence under a shared ontological interpretation.


%The detection of PT/DT divergence is commonly addressed through runtime monitoring and trace analysis. Munoz et al.~\cite{10.1145/3652620.3688267} detect divergence by aligning observed PT traces with DT behavior and applying anomaly detection techniques, producing verdicts that depend on the executions observed and the properties instrumented. Similarly, Moon et al.~\cite{moon2023consistency} check PT/DT consistency by verifying that observed timed event traces conform to a DT model expressed as a Time Colored Petri Net. Compared to our approach, they rely on syntactic correspondence between traces or model executions, whereas our framework grounds PT and DT observations in a shared ontology, enabling behavioral equivalence to be established semantically even when their observable vocabularies differ. Similarly to our framework, ontologies are used in~\cite{kamburjan2022twinningbyconstruction,Gil2025} to lift DT runtime states into knowledge graphs, enabling semantic querying and monitoring of PT/DT correspondence. 

%Generally, while run-time monitoring aim to establish alignment between the PT and DT, they do so operationally and over selected constraints, whereas we establish consistency completely at design time, over all reachable states, and for all properties expressible over the aligned fragment. Beyond the stronger guarantee, our approach eliminates the engineering overhead of implementing, deploying, and maintaining runtime consistency mechanisms, avoiding as a result a burden that compounds with system evolution since every change requires updating and re-validating the monitoring infrastructure. Our refinement theory handles ontology evolution directly, ensuring that deliberate updates to domain knowledge do not invalidate previously established alignment without requiring any monitor to be rewritten or redeployed. Runtime monitoring nevertheless remains complementary for detecting deployment-time deviations such~\cite{10.1145/3735948.3736156} or quantify uncertainty in DT predictions under incomplete environment models~\cite{Woodcock2021}, which are phenomena outside the scope of any design-time framework. Our contribution is therefore not to replace runtime monitoring, but to discharge by construction the semantic consistency problem that existing approaches address operationally.

%Another line of work studies the co-evolution of PT and DT models. Model-driven toolchains propagate structural changes from the PT to maintain DT infrastructure~\cite{lehner2021aml4dt,Mertens2024}. These approaches ensure structural correspondence between twins but do not address semantic or behavioral equivalence. Our framework is complementary: it does not prescribe how the DT evolves, but provides formal guarantees that the resulting models remain semantically aligned with respect to domain knowledge.


%Recent work has explored formal specification and verification of digital twin behavior. Huang et al.~\cite{huang2024formaldt} employ the Temporal Logic of Actions (TLA) to specify synchronization and information-flow properties of DT architectures, enabling formal verification of correct DT operation. Liu et al.~\cite{liu2023spn} introduce Semantic Petri Nets, which integrate RDF/OWL knowledge bases with Petri Net models to validate temporal rules governing DT state transitions through SPARQL-based execution. While these approaches bring formal reasoning into DT modeling, their focus remains on verifying properties of the digital twin itself or validating rule-based constraints during model execution. In contrast, our work formalizes the semantic alignment between PT and DT as a behavioral equivalence relation under a shared ontological interpretation. Rather than checking selected rules or properties over DT executions, our framework establishes design-time guarantees that PT and DT behaviors are semantically consistent across all reachable states and for all properties expressible over the aligned fragment.